\documentclass[12pt,a4paper]{article}

\usepackage{fontawesome5}
\providecommand{\faPencilAlt}{\faIcon{pencil-alt}}

\usepackage{styles/couleurs}
\usepackage{styles/mise_en_page}
\usepackage{styles/boites}
\usepackage{styles/maths}
\usepackage{styles/figures}

\usepackage{enumitem}
\usepackage{array}
\usepackage{tabularx}
\usepackage{amsmath}
\usepackage{amssymb}
\usepackage{tikz}
\usepackage{xparse}

\usetikzlibrary{calc,angles,quotes,arrows.meta}

\setlist[itemize]{leftmargin=1.4em,itemsep=0.15em,topsep=0.2em}
\setlist[enumerate]{leftmargin=*,itemsep=0.22em,topsep=0.25em}

\newcommand{\lignepointillee}{\noindent\dotfill\par\vspace{0.45em}}

\newcommand{\lignesreponse}[1]{%
\begin{tcolorbox}[
enhanced,colback=white,colframe=NeutreC!40,boxrule=0.5pt,arc=1.5mm,
left=3mm,right=3mm,top=2mm,bottom=1mm,before skip=2pt,after skip=2pt]
\foreach \i in {1,...,#1}{\lignepointillee}
\end{tcolorbox}}

\newcommand{\titreseq}[1]{%
\begin{center}
{\sffamily\LARGE\bfseries\color{TitreC}Séquence 6 -- Cosinus d'un angle aigu}\\[0.2em]
{\sffamily\large\bfseries\color{SavoirC}#1}
\end{center}}

\begin{document}

\titreseq{Activité 4 -- Calculer le côté adjacent}

\section*{Partie A -- Comprendre l'opération}

Dans le triangle \(ABC\) rectangle en \(B\) :

\[
AC=12\text{ cm}
\qquad\text{et}\qquad
\widehat{BAC}=37^\circ.
\]

On cherche \(AB\).

\begin{center}
\begin{tikzpicture}[scale=0.85]
  \coordinate (A) at (0,0);
  \coordinate (B) at (5.2,0);
  \coordinate (C) at (5.2,3.9);

  \draw[line width=0.9pt] (A)--(B)--(C)--cycle;
  \draw (4.82,0)--(4.82,0.38)--(5.2,0.38);

  \node[left] at (A) {$A$};
  \node[below] at (B) {$B$};
  \node[above] at (C) {$C$};

  \node[above left] at (2.7,2.2) {$12$ cm};
  \node[below] at (2.7,0) {$AB\,?$};

  \draw[ActionC,line width=1pt] (0.8,0) arc[start angle=0,end angle=37,radius=0.8];
  \node[ActionC] at (1.15,0.43) {$37^\circ$};
\end{tikzpicture}
\end{center}

Complète.

\[
\cos(37^\circ)=\frac{\ldots}{\ldots}
\]

Le cosinus est le nombre qui multiplie l'hypoténuse pour obtenir le côté adjacent :

\[
AB=AC\times\cos(37^\circ).
\]

Donc :

\[
AB=12\times\cos(37^\circ)\approx\ldots\text{ cm}.
\]

Le résultat est-il cohérent avec \(AB<AC\) ?

\section*{Partie B -- Méthode de rédaction}

Complète le modèle.

\begin{tcolorbox}[colback=white,colframe=SavoirC,boxrule=0.7pt,arc=2mm]
Dans le triangle \(ABC\) rectangle en \(B\), on a :

\[
\cos(\widehat{BAC})=\frac{AB}{AC}.
\]

Donc :

\[
AB=AC\times\cos(\widehat{BAC}).
\]

Ainsi :

\[
AB=12\times\cos(37^\circ)
\approx\ldots\text{ cm}.
\]
\end{tcolorbox}

\section*{Partie C -- S'entraîner}

Pour chaque figure :

\begin{enumerate}
  \item nomme le triangle rectangle ;
  \item écris le cosinus de l'angle donné ;
  \item calcule la longueur demandée ;
  \item arrondis au dixième et contrôle le résultat.
\end{enumerate}

\begin{minipage}[t]{0.31\textwidth}
\centering
\begin{tikzpicture}[scale=0.55]
  \coordinate (J) at (0,0);
  \coordinate (V) at (4.5,0);
  \coordinate (K) at (4.5,2.2);
  \draw[line width=0.9pt] (J)--(V)--(K)--cycle;
  \draw (4.15,0)--(4.15,0.35)--(4.5,0.35);
  \node[left] at (J) {$J$};
  \node[below] at (V) {$V$};
  \node[above] at (K) {$K$};
  \node[above left] at (2.3,1.25) {$8$ cm};
  \node[below] at (2.4,0) {$JV\,?$};
  \draw[ActionC] (0.75,0) arc[start angle=0,end angle=26,radius=0.75];
  \node[ActionC] at (1.05,0.35) {$20^\circ$};
\end{tikzpicture}
\end{minipage}
\hfill
\begin{minipage}[t]{0.31\textwidth}
\centering
\begin{tikzpicture}[scale=0.55]
  \coordinate (R) at (0,0);
  \coordinate (S) at (4.4,0);
  \coordinate (T) at (4.4,3.1);
  \draw[line width=0.9pt] (R)--(S)--(T)--cycle;
  \draw (4.05,0)--(4.05,0.35)--(4.4,0.35);
  \node[left] at (R) {$R$};
  \node[below] at (S) {$S$};
  \node[above] at (T) {$T$};
  \node[above left] at (2.3,1.65) {$15$ cm};
  \node[below] at (2.3,0) {$RS\,?$};
  \draw[ActionC] (0.75,0) arc[start angle=0,end angle=35,radius=0.75];
  \node[ActionC] at (1.05,0.42) {$35^\circ$};
\end{tikzpicture}
\end{minipage}
\hfill
\begin{minipage}[t]{0.31\textwidth}
\centering
\begin{tikzpicture}[scale=0.55]
  \coordinate (D) at (0,0);
  \coordinate (E) at (4.2,0);
  \coordinate (F) at (4.2,4.3);
  \draw[line width=0.9pt] (D)--(E)--(F)--cycle;
  \draw (3.85,0)--(3.85,0.35)--(4.2,0.35);
  \node[left] at (D) {$D$};
  \node[below] at (E) {$E$};
  \node[above] at (F) {$F$};
  \node[above left] at (2.2,2.3) {$9{,}5$ cm};
  \node[below] at (2.2,0) {$DE\,?$};
  \draw[ActionC] (0.75,0) arc[start angle=0,end angle=46,radius=0.75];
  \node[ActionC] at (1.02,0.5) {$46^\circ$};
\end{tikzpicture}
\end{minipage}

\lignesreponse{10}

\aretenir{Pour calculer le côté adjacent :

\[
\boxed{\text{côté adjacent}
=
\text{hypoténuse}\times\cos(\text{angle})}
\]

Comme le cosinus d'un angle aigu est inférieur à \(1\),
le résultat doit être plus petit que l'hypoténuse.}

\end{document}
